
\section{Sense-Labelled Evaluation}\label{app:sense-labelled-evaluation}

\S\ref{subsec:same-token-analysis} reads individual scores for a fixed
token under two contexts. This appendix tests the same feature groups
against labelled sense data and measures what the sparse account gives
up against the full decoded state.

\paragraph{Data and coverage.}
CoarseWSD-20 \citep{loureiro2021coarsewsd} supplies 20 ambiguous words
with coarse sense labels and a fixed train and test split, 33{,}566
contexts in total. We keep the
shipped splits and score the final position of a cloze prompt that asks
for the ambiguous word. All 20 words resolve to a single token under
the Qwen3.5 vocabularies. Seven do not under the R1-Llama-8B
vocabulary, so that model is evaluated on 13 words and 16{,}785
contexts and its column is not comparable item for item with the other
two. Row-gate coverage is 1.00 on the Qwen models and 0.976 on
R1-Llama-8B. Sense distributions are skewed, so we report balanced
accuracy throughout. Unbalanced accuracy gives the same ordering and is
released with the run artifacts.

\paragraph{Protocol.}
For each word we select one group of eight features per sense on the
training split by largest mean contribution difference between that
sense and the others, then assign each held-out context to the sense
whose group carries the larger summed contribution. Feature scales are
standardized by training statistics alone, and no test label enters the
selection. The null permutes sense labels within the training split and
repeats the whole selection, so it inherits the selection procedure.

\paragraph{Sense alignment.}
\Cref{tab:app-sense-alignment} gives the result. Selected groups recover
held-out senses at 0.90, 0.92, and 0.80 balanced accuracy against a
permutation null near 0.41. Every word on Qwen3.5-2B, 19 of 20 on
Qwen3.5-9B, and 12 of 13 on R1-Llama-8B exceed that null's 95th
percentile, and every word on all three models exceeds its majority
baseline. The finding does not depend on the account size. Across
sizes one, two, four, and eight the majority baseline is exceeded for
every word on every model, and the count above the null's 95th
percentile moves by at most one word.

The last two rows of the same table bound the result from above. A
nearest centroid probe on the full decoded state reaches 0.96, 0.97, and
0.92, and the full 256-dimensional account without any selection
reaches 0.95, 0.97, and 0.89. The three rows form a sparsity
ladder in which factorizing the row costs little. Compressing to eight features
costs a further 5 to 9 points against the full account, and 5 to 12
points against the full state. That is what the account costs, and it
buys a set of features a reader can name.

This is an alignment result for the selected groups. The comparison
against direct \(\WU\) geometry is made on reconstruction and coverage
in \Cref{tab:readout-score-fidelity-summary} and on grouping recovery in
\S\ref{subsec:stability}, where the alternatives are matched by
construction.

\paragraph{A classifier framing.}
A different framing trains a per-word classifier on the ten features
that carry the largest absolute contribution, skipping the group
selection. Under that framing the unweighted projections
\(p_i(\rstate)\) beat the signed contributions on all three models, at
0.87 against 0.84, 0.90 against 0.76, and 0.73 against 0.67. Row-coefficient shuffles and randomly drawn features sit close to
the signed contributions.

That ordering follows from the factorization. Holding a vocabulary row \(v\) fixed
freezes its coefficients \(z_{v,i}\) and leaves only \(p_i(\rstate)\) free
(\S\ref{sec:method-local-score}), so for a fixed word every weighted
variant is a constant per-coordinate rescaling of the same projections.
A per-word supervised classifier gains nothing from that rescaling, and
a shuffle of the coefficients preserves the projections it is meant to
remove. The signed contributions name which readout directions support
a selected score. They are not a sense representation, and this second
framing measures the second quantity. Sense structure in embedding
geometry has been studied directly \citep{arora2018polysemy}.

\paragraph{Scope.}
The evaluation covers one dataset, one sparsity setting per model, and
one prompt template. No behavioral filter is applied, and the model
ranks the ambiguous word at median rank 174, 10, and 1131 on the three
models, so a share of the scored contexts sits outside the operating
regime of Appendix~\ref{subsec:interpretation-boundary}. Filtering on
target rank would change the reported values in an untested direction.

\begin{table}[!tbp]
\caption{\textbf{Sense alignment on CoarseWSD-20.} Balanced accuracy on
held-out contexts for groups of eight features selected on the training
split, against a label-permutation null that repeats the selection, a
majority baseline, and two references that use no selection.}
\label{tab:app-sense-alignment}
\centering
\scriptsize
\setlength{\tabcolsep}{4.0pt}
\renewcommand{\arraystretch}{1.08}
\begin{tabularx}{\linewidth}{@{}>{\RaggedRight\arraybackslash}Xccc@{}}
\toprule
 & Qwen3.5-2B & Qwen3.5-9B & R1-Llama-8B \\
Words & 20 & 20 & 13 \\
\midrule
Selected groups & 0.899 & 0.923 & 0.798 \\
Permutation null & 0.408 & 0.407 & 0.422 \\
Majority & 0.410 & 0.410 & 0.429 \\
Words above null p95 & 20/20 & 19/20 & 12/13 \\
Words above majority & 20/20 & 20/20 & 13/13 \\
\midrule
Full account, no selection & 0.947 & 0.970 & 0.892 \\
Hidden state & 0.958 & 0.971 & 0.920 \\
\bottomrule
\end{tabularx}
\end{table}
